\documentclass[11pt]{article}

% --- arXiv-safe preamble: only packages present in the arXiv TeX Live tree ---
\usepackage[utf8]{inputenc}
\usepackage[T1]{fontenc}
\usepackage{lmodern}
\usepackage[margin=1in]{geometry}
\usepackage{amsmath}
\usepackage{booktabs}
\usepackage{array}
\usepackage{url}
\usepackage[hidelinks]{hyperref}
\usepackage{microtype}

\title{Measuring the Model Context Protocol Server Ecosystem:\\
A Population-Scale Census and Conformance Study}

\author{Haseeb Mohammed Afsar\\
  \small Independent researcher\\
  \small ORCID: \href{https://orcid.org/0009-0000-4038-1272}{0009-0000-4038-1272}}

\date{July 2026}

\begin{document}
\maketitle

\begin{abstract}
The Model Context Protocol (MCP) is an open standard that lets large language
model (LLM) applications discover and call external tools through a uniform
JSON-RPC~2.0 interface. Its adoption has produced a large, rapidly growing
population of independently authored servers, but the ecosystem has not been
measured: basic questions about its size, how servers ship, which transports and
spec revisions are in use, and whether servers conform to the specification and
behave reliably have gone unanswered. We report a two-tier empirical study. A
\emph{census tier} sweeps the entire official MCP registry (16{,}548 unique servers
at a July~2026 snapshot) and measures only self-declared metadata, so it runs no
third-party code and scales to the whole published population. A \emph{dynamic tier}
launches a credential-free, stdio-launchable subset and measures actual behavior with
the open-source \texttt{mcp-probe} tool: handshake success, JSON~Schema validity of
every advertised tool, presence of safety annotations, protocol-version negotiation,
and latency. Three results stand out. First, the population is dominated by hosted
remote servers (42.6\% remote-only) and \texttt{npm}/\texttt{stdio} local packages,
and it has already fragmented across five \texttt{\$schema} revisions. Second,
\emph{hard} conformance is high: all 200 tools measured in the behavioral sample carry
valid JSON~Schema \texttt{inputSchema} definitions with zero fatal violations. Third,
the real variance is in optional \emph{safety} metadata: 41.5\% of tools omit the
\texttt{annotations} (\texttt{readOnlyHint}, \texttt{destructiveHint}, \dots) that
tell an autonomous agent whether a tool is safe to call \emph{before} calling it, and
the omission is a clean server-level, all-or-nothing decision rather than per-tool
oversight. We release the full dataset and a re-runnable pipeline so the measurement
can be repeated at each spec release as a longitudinal series.
\end{abstract}

\section{Introduction}

The Model Context Protocol (MCP)~\cite{mcp2024} is an open standard for connecting
LLM applications to external tools, resources, and prompts over a uniform JSON-RPC
2.0 interface. Since its introduction the number of independently authored MCP
servers has grown quickly, and with it the variance in their quality: tool
\texttt{inputSchema} definitions can drift out of spec, handshakes can fail on edge
cases, and malformed responses may surface only when an LLM constructs an invalid
tool call at run time. As agentic systems move into production, the reliability of
the tools they call becomes a first-order concern~\cite{toolreliability}.

Despite this, the MCP server ecosystem has not been measured empirically. We do not
know how large the published population is, how servers are actually shipped and
consumed (hosted remote endpoints versus locally installed packages), which
transports and specification revisions are in use, or whether servers conform to the
specification and behave reliably in practice. This paper answers those questions
with a reproducible, two-tier study and releases the underlying dataset and tooling.

\paragraph{Contributions.}
\begin{itemize}
  \item A \textbf{population-scale census} of the official MCP registry
    (16{,}548 unique servers), classified by deployment model, package ecosystem,
    transport, lifecycle status, and \texttt{\$schema} revision
    (Section~\ref{sec:census}). Because it measures only self-declared metadata it
    runs no third-party code and covers the whole published population.
  \item A \textbf{behavioral conformance study} of a credential-free,
    stdio-launchable sample using the open-source \texttt{mcp-probe} instrument,
    reporting schema validity, safety-annotation coverage, protocol-version drift,
    and latency (Section~\ref{sec:dynamic}).
  \item A \textbf{released, re-runnable pipeline and dataset}
    (Section~\ref{sec:data}) so the measurement can be repeated at each spec release,
    producing a longitudinal series rather than a one-off snapshot.
\end{itemize}

\section{Study design}
\label{sec:design}

\subsection{Research questions}
\begin{description}
  \item[RQ0 (census).] How large is the published MCP server population, and how is
    it distributed across deployment models, package ecosystems, transports, and
    specification revisions?
  \item[RQ1 (conformance).] What fraction of public MCP servers initialize
    successfully and advertise spec-valid tool \texttt{inputSchema} definitions?
  \item[RQ2 (drift).] Which schema/spec deviations are most common, and do they
    cluster by authoring pattern?
  \item[RQ3 (reliability).] What is the latency distribution of protocol operations?
\end{description}

\subsection{Two measurement tiers}
The questions above have different \emph{safe} sample sizes, so we measure at two
scales.

\paragraph{Census tier (population, code-free).}
A harvester sweeps the entire official MCP registry
(\url{https://registry.modelcontextprotocol.io/v0/servers}) via cursor pagination
and records only metadata the registry publishes: deployment model, package
ecosystem, declared transport, lifecycle status, and pinned \texttt{\$schema}
revision. Because no server code is executed, this tier safely scales to the whole
published population and forms the denominator the dynamic tier samples from. The
sweep records page count, version-row count, and a \texttt{sweepComplete} flag; if a
pagination backstop is ever hit the run is marked a lower bound rather than silently
truncated.

\paragraph{Dynamic tier (sampled, behavioral).}
A credential-free, stdio-launchable subset is launched over stdio and probed with
\texttt{mcp-probe}, a dependency-light command-line tool that speaks the MCP wire
protocol directly, performs the \texttt{initialize} handshake, enumerates advertised
tools via \texttt{tools/list}, validates each tool's JSON~Schema~\cite{jsonschema}
against the constraints the MCP specification places on tool definitions, and times
protocol operations. Untrusted code executes in this tier, so it is a vetted,
sampled subset rather than the whole population.

\subsection{Sample frame}
The census frame is every server in the official registry at a snapshot timestamp
recorded with the dataset. The dynamic frame is the reference set
(\texttt{modelcontextprotocol/servers}) plus the \texttt{npm}-published,
stdio-launchable, active servers the census emits as candidates. Credential needs are
verified empirically at run time; every included and excluded server is recorded with
a reason so the frame is auditable.

\subsection{Ethics and guardrails}
Public servers only. No credentials are ever supplied. Only read-only or explicitly
safe calls are exercised; no side-effecting tool is invoked. Findings are reported in
aggregate as ecosystem health, not as a leaderboard of individual failures.

\section{RQ0: population census}
\label{sec:census}

At the July~2026 snapshot the sweep fetched 51{,}151 version rows across 512 registry
pages and resolved 16{,}548 unique servers (latest version each), of which 16{,}377
are marked active. For context, GitHub reports 20{,}629 repositories under the
\texttt{topic:mcp-server} tag; we report this only as an indication of how much
larger the informal universe is than the registry frame we measure.

\paragraph{Deployment model.}
Table~\ref{tab:deploy} shows how servers ship. A 42.6\% remote-only share means most
servers are hosted HTTP/SSE endpoints rather than local stdio tools: the ecosystem is
consumed as much through network integrations (with the attendant auth and
reliability concerns) as through locally installed packages.

\begin{table}[h]
\centering
\begin{tabular}{lrr}
\toprule
Deployment model & Count & Share \\
\midrule
Package-only (installed locally) & 8340 & 50.4\% \\
Remote-only (hosted HTTP/SSE)    & 7057 & 42.6\% \\
Both                             & 852  & 5.1\%  \\
Neither declared                & 299  & 1.8\%  \\
\bottomrule
\end{tabular}
\caption{Deployment model across the 16{,}548-server population.}
\label{tab:deploy}
\end{table}

\paragraph{Package ecosystems and transports.}
Of servers that ship a package, \texttt{npm} dominates (5880; 35.5\% of all servers),
followed by \texttt{pypi} (2648; 16\%), \texttt{oci} (562), \texttt{mcpb} (349),
\texttt{nuget} (83), and a single \texttt{cargo} crate. On transports, package servers
overwhelmingly declare \texttt{stdio} (9057), while remote servers overwhelmingly
declare \texttt{streamable-http} (7478) with a long tail of \texttt{sse} (682) ---
consistent with the spec's shift toward streamable HTTP for hosted deployments.

\paragraph{Specification-version drift.}
Registry entries pin a \texttt{\$schema} revision, giving a clean, code-free drift
signal across the whole population (Table~\ref{tab:schema}). A single revision
(\texttt{2025-12-11}) accounts for 88.2\% of servers, but five revisions are live
simultaneously, and 11.8\% of the population trails the current one --- a
population-scale confirmation that the specification is a moving target the ecosystem
lags.

\begin{table}[h]
\centering
\begin{tabular}{lrr}
\toprule
\texttt{\$schema} revision & Count & Share \\
\midrule
2025-12-11 & 14603 & 88.2\% \\
2025-09-29 & 979   & 5.9\%  \\
2025-10-17 & 556   & 3.4\%  \\
2025-07-09 & 271   & 1.6\%  \\
2025-09-16 & 137   & 0.8\%  \\
unknown    & 2     & 0.0\%  \\
\bottomrule
\end{tabular}
\caption{Schema-revision spread across the population.}
\label{tab:schema}
\end{table}

\paragraph{Dynamic-tier denominator.}
5804 servers (35.1\% of the population) are \texttt{npm}-published stdio servers and
thus, in principle, launchable by \texttt{npx} for behavioral probing. This is the
frame the dynamic tier samples from; actual inclusion is decided empirically because
many require credentials or a mandatory launch argument.

\section{RQ1--RQ3: behavioral conformance}
\label{sec:dynamic}

We report a baseline dynamic run over a reproducible sample of 24 candidate servers
(the official reference servers plus a set of popular community servers), each
launched over stdio with no credentials. Sixteen initialized without credentials and
were included (200 tools total); eight were excluded and recorded with a reason (four
require credentials; four fail the handshake without a required argument such as a
connection string or target URL). No server was dropped silently.

\subsection{RQ1: conformance}
All 16 included servers completed the \texttt{initialize} handshake and returned a
tool list. All \textbf{200 of 200 (100\%)} tool \texttt{inputSchema} definitions are
valid JSON~Schema with \textbf{zero fatal violations}: no missing schemas, no invalid
\texttt{type}, no malformed \texttt{properties}/\texttt{required}. 117 of 200 tools
(58.5\%) are fully clean --- valid schema \emph{and} no warning-level issues. At the
level of hard schema validity, the surveyed servers are strongly conformant; the
intuition that MCP tool schemas are frequently malformed is \emph{not} supported. This
held as the frame nearly doubled from a nine-server pilot (104 to 200 tools, 100\% to
100\%), which suggests it is a structural property rather than a small-sample artifact.

\subsection{RQ2: drift}
The dominant deviation is a single one: \emph{missing tool \texttt{annotations}}
(the \texttt{readOnlyHint} / \texttt{destructiveHint} / \texttt{idempotentHint} /
\texttt{openWorldHint} metadata introduced in a later spec revision).
\textbf{83 of 200 tools (41.5\%)} omit it. Critically, the gap is \textbf{strictly
server-level, not random}: the split is all-or-nothing with \emph{zero} partial
servers --- 8 servers annotate every tool and 8 annotate none, with not a single
server in between. Protocol-version fragmentation appears in the behavioral sample
too: two versions are live (\texttt{2025-06-18} on 12 servers and the older
\texttt{2024-11-05} on 4), and one official reference server is simultaneously on the
old protocol version \emph{and} annotates no tools.

Because annotations are how a client or agent learns whether a tool is read-only or
potentially destructive \emph{before} invoking it, a 41.5\% omission rate means a
meaningful fraction of tools give an autonomous agent no machine-readable signal about
their side effects. That the omission is a clean server-level split (never partial)
indicates a per-project authoring decision rather than per-tool oversight --- a
measurable, actionable gap between the specification's safety affordances and
ecosystem adoption.

\subsection{RQ3: reliability}
Handshake round-trip across the 16 servers was p50~$\approx$~1.13~s,
p95~$\approx$~10.41~s (min 0.64~s, max 10.41~s). We flag this as a \emph{contaminated}
measurement: it includes first-run \texttt{npx} install cost. A hardened latency run
must pre-install and warm servers first; we report the baseline here only to fix the
method, not as a clean latency estimate.

\section{Discussion}

The two tiers tell a consistent story. At population scale the ecosystem is large,
majority-hosted, \texttt{npm}/\texttt{stdio}-led where local, and already fragmented
across five spec revisions with an 11.8\% trailing tail. At the behavioral level, hard
conformance is essentially solved --- schema validity is 100\% --- but the safety
layer the spec added on top (annotations) has been adopted by only some authors, in an
all-or-nothing manner. For agent builders, the practical implication is that schema
\emph{shape} can be trusted but tool \emph{safety metadata} cannot be assumed present;
clients that reason about destructiveness before calling a tool must treat missing
annotations as the common case, not the exception.

\section{Threats to validity}
The census measures self-declared registry metadata; a server misdeclaring its
transport or schema is counted as declared. The registry, while authoritative, is one
view of the ecosystem: servers never published to it are out of frame (hence the
GitHub topic count reported only as context). The dynamic sample is small
($n=16$ servers) and restricted to GitHub-visible, \texttt{npm}-launchable,
credential-free servers, so its absolute percentages are expected to move as the frame
grows; the reusable contribution is the method and the categories of deviation it
surfaces. Latency is contaminated by install cost as noted. All numbers are tagged to
the snapshot date and the client spec baseline (\texttt{2025-06-18}); the pipeline is
re-runnable to produce a longitudinal series.

\section{Data and code availability}
\label{sec:data}
The \texttt{mcp-probe} tool, the harvester, the study protocol, and the full raw
per-server records are open source at
\url{https://github.com/itguruhaseeb/mcp-probe}. The versioned dataset and tool
release are archived on Zenodo under the concept DOI
\href{https://doi.org/10.5281/zenodo.21347997}{10.5281/zenodo.21347997}, which always
resolves to the latest version~\cite{mcpprobe_dataset}. Re-running
\texttt{node benchmark/harvest.mjs} reproduces the census, and
\texttt{node benchmark/run-study.mjs} reproduces the behavioral run.

\section*{Acknowledgements}
We thank the maintainers of the Model Context Protocol specification and its
reference server implementations, against which \texttt{mcp-probe} is tested.

\begin{thebibliography}{9}

\bibitem{mcp2024}
Anthropic.
\newblock \emph{Model Context Protocol Specification}.
\newblock Open standard for connecting large language model applications to external
  tools and data, 2024.
\newblock \url{https://modelcontextprotocol.io}.

\bibitem{jsonschema}
A.~Wright, H.~Andrews, B.~Hutton, and G.~Dennis.
\newblock \emph{JSON Schema: A Media Type for Describing JSON Documents}.
\newblock IETF Internet-Draft, 2022.
\newblock \url{https://json-schema.org/specification}.

\bibitem{mcpprobe_dataset}
H.~Mohammed Afsar.
\newblock \emph{mcp-probe: a conformance and reliability checker for Model Context
  Protocol servers} (software and dataset).
\newblock Zenodo, 2026.
\newblock \url{https://doi.org/10.5281/zenodo.21347997}.

\bibitem{toolreliability}
S.~Yao, N.~Shinn, P.~Razavi, and K.~Narasimhan.
\newblock \emph{$\tau$-bench: A Benchmark for Tool-Agent-User Interaction in
  Real-World Domains}.
\newblock arXiv:2406.12045, 2024.
\newblock \url{https://doi.org/10.48550/arXiv.2406.12045}.

\end{thebibliography}

\end{document}
